Surface area is one of the most requested measurements in X-ray computed tomography (CT) of materials
and parts. It controls reaction rates in catalysts, flow in porous media and the quality of additively
manufactured lattices. It is usually measured from the reconstructed voxel volume, either by counting
the exposed faces of segmented voxels or by extracting a triangle mesh with marching cubes
\citep{lorensen1987} and summing the triangle areas.

Both routes are biased. Counting voxel faces measures the area of a staircase, which is about 50\%
larger than the true surface at every resolution. Marching cubes on a thresholded (binary) volume is
better, but its error does not vanish as the voxels get smaller \citep{kenmochi2000,svane2013}.
Marching cubes on the greyscale volume is accurate for clean, well-resolved data, but real CT volumes
are noisy and many structures are only a few voxels across. Because the error depends on resolution,
noise, curvature and thin features at the same time, a single correction factor cannot remove it.

This work builds a synthetic CT-like dataset with exactly known surface areas and trains models that
predict a correction to the greyscale marching-cubes area for each object, so that learning acts as a
correction layer on top of the classical measurement rather than replacing it.
